% !PRINTED 303
\weberpage{303}

%<<CONT
also
\[
q=p,\qquad p=3
\]
sein, was wir ausgeschlossen haben.

Im Falle 2. aber folgt für jedes beliebige $p$:

\[ q=2. \]

Daher muß in $\mathfrak{K}$ jedenfalls die Substitution $B$ vorkommen, und die gesuchte Gruppe ist durch (13), (14), (16), (17) erschöpft.

Wir zeigen zunächst, daß in den beiden Substitutionen $V, V'$ der Gruppe $\mathfrak{K}$ keine der Zahlen $\alpha,\beta,\gamma,\delta,\alpha',\beta',\gamma',\delta'$ kongruent $0$ sein kann.

Es kommt nämlich, wie wir schon gesehen haben, in $\mathfrak{K}$ nicht vor

\[ \begin{pmatrix}1, & 0\\ \gamma, & 1\end{pmatrix}, \]

wenn $\gamma$ von $0$ verschieden ist, also auch nicht

\[ \begin{pmatrix}\alpha, & 0\\ 0, & \alpha^{-1}\end{pmatrix}\begin{pmatrix}1, & 0\\ \alpha\gamma, & 1\end{pmatrix}=\begin{pmatrix}\alpha, & 0\\ \gamma, & \alpha^{-1}\end{pmatrix}, \]

wenn $\alpha$ beliebig ist; und folglich auch nicht

\begin{align*}
\begin{pmatrix}\alpha, & 0\\ \gamma, & \alpha^{-1}\end{pmatrix}\begin{pmatrix}0, & 1\\ -1, & 0\end{pmatrix} &=\begin{pmatrix}0, & \alpha\\ -\alpha^{-1}, & \gamma\end{pmatrix}\\
\begin{pmatrix}0, & 1\\ -1, & 0\end{pmatrix}\begin{pmatrix}\alpha, & 0\\ \gamma, & \alpha^{-1}\end{pmatrix} &=\begin{pmatrix}\gamma, & \alpha^{-1}\\ -\alpha, & 0\end{pmatrix}\\
\begin{pmatrix}0, & -1\\ 1, & 0\end{pmatrix}\begin{pmatrix}0, & \alpha\\ -\alpha^{-1}, & -\gamma\end{pmatrix} &=\begin{pmatrix}\alpha^{-1}, & \gamma\\ 0, & \alpha\end{pmatrix},
\end{align*}

worin alle Substitutionen $\begin{pmatrix}\alpha, & \beta\\ \gamma, & \delta\end{pmatrix}$ enthalten sind, in denen eine der vier Zahlen $\alpha,\beta,\gamma,\delta$ kongruent mit $0$ ist.

Die Substitutionen des Systems (16) bezeichnen wir jetzt mit

\[ W=\begin{pmatrix}\alpha g^{r+s}, & \beta g^{r-s}\\ \gamma g^{-r+s}, & \delta g^{-r-s}\end{pmatrix}, \]

und setzen

\weq{18}{\xi\equiv\alpha g^{r+s},\qquad \varrho\equiv\alpha^{-1}\beta g^{-2s}}

\weq{19}{\alpha\delta\equiv m,\qquad \beta\gamma\equiv m-1\quad(\mathrm{mod}\ p).}

$m$ ändert sich nicht, wenn $V$ durch eine beliebige Substitution des Systems (16) ersetzt wird. Die Zahl $m'$ soll die entsprechende Bedeutung für das System (17) haben; $\xi$ kann in jedem dieser beiden Systeme jeden der Werte

\weq{20}{1,2,\dots p-1}
%CONT>>
